\chapter{Pomijanie typów, modele atomowe i pierwsze}
\label{ch:mod-omitting-types}
\section{Motywacja i typy niegłówne}
W rozdziale \ref{ch:mod-types} poznaliśmy przestrzenie
typów, a w rozdziale \ref{ch:mod-saturation} badaliśmy
modele realizujące wiele typów. Teraz stawiamy pytanie
odwrotne: czy można zbudować model, w którym żaden
element (lub żadna skończona krotka) nie realizuje
wskazanego typu?

Pracujemy w ZFC, w logice pierwszego rzędu z równością,
z niepustymi uniwersami. Wszystkie typy w tym rozdziale
mają skończoną dodatnią liczbę zmiennych i, o ile nie
zaznaczono inaczej, nie mają parametrów. Typ częściowy
nad spełnialną teorią $T$ oraz jego realizację
i pomijanie rozumiemy jak w definicji
\ref{def:mod-partial-type}. Typ zupełny nad $T$
jest punktem $S_n(T)$ (definicja
\ref{def:mod-complete-type}).

\begin{definicja}[Typ główny i typ niegłówny]
\label{def:mod-ott-principal}
Niech $p\subseteq F_n(\mathcal L)$ będzie typem
częściowym nad spełnialną teorią $T$.
Nazywamy go głównym, jeżeli istnieje formuła
$\theta(\bar x)$ realizowalna w pewnym modelu $T$
i taka, że dla każdej $\varphi\in p$
\[
T\models\forall\bar x\,(\theta(\bar x)\to\varphi(\bar x)).
\]
Jeżeli nie istnieje taka formuła, typ nazywamy
niegłównym. Formuła $\theta$ nie musi należeć do $p$.
\end{definicja}

Warunek realizowalności $\theta$ jest niezbędny:
formuła sprzeczna implikuje semantycznie każdą
formułę, ale nie opisuje żadnej krotki.
Dla typów częściowych główność w tym znaczeniu
nie oznacza, że typ jest równoważny jednej formule:
wystarczy, aby pewna niepusta część przestrzeni
typów była zawarta w zbiorze jego realizacji.

\begin{lemat}[Główność a topologia przestrzeni typów]
\label{lem:mod-ott-interior}
Dla typu częściowego $p$ nad $T$ połóżmy
\[
C_p=\{q\in S_n(T):p\subseteq q\}
=\bigcap_{\varphi\in p}[\varphi].
\]
Zbiór $C_p$ jest niepusty i domknięty.
Typ $p$ jest główny dokładnie wtedy, gdy $C_p$
ma niepuste wnętrze. W szczególności typ zupełny
jest główny dokładnie wtedy, gdy jest izolowany.
\end{lemat}
\begin{proof}
Niepustość wynika z twierdzenia
\ref{thm:mod-type-completion}. Każdy zbiór
$[\varphi]$ jest domknięty i otwarty w topologii
Stone'a (lemat \ref{lem:mod-stone-basic}), więc
ich część wspólna jest domknięta.
Jeżeli $\theta$ świadczy o główności $p$, to
$[\theta]$ jest niepustym zbiorem otwartym i
$[\theta]\subseteq C_p$, ponieważ każda realizacja
typu zawierającego $\theta$ spełnia wszystkie
formuły $p$. Odwrotnie, jeśli $C_p$ ma niepuste
wnętrze, wybieramy niepusty zbiór bazowy
$[\theta]\subseteq C_p$. Realizowalność $\theta$
wynika z niepustości $[\theta]$.
Gdyby dla pewnej $\varphi\in p$ implikacja
$\theta\to\varphi$ nie zachodziła modulo $T$,
krotka realizująca $\theta\land\neg\varphi$
w modelu $T$ miałaby typ należący do
$[\theta]\setminus C_p$, co jest sprzecznością.
Dla typu zupełnego $p$ maksymalność daje
$C_p=\{p\}$, więc niepuste wnętrze jest
równoważne izolacji punktu w sensie definicji
\ref{def:mod-isolated-type}.
\end{proof}

W szczególności zbiór $C_p$ dla typu niegłównego
jest domknięty i ma puste wnętrze, czyli jest
nigdzie gęsty. Intuicyjnie: każde realizowalne
skończone wymaganie można wzmocnić tak, aby
przynajmniej jedna formuła niegłównego typu
była fałszywa.

\begin{lemat}[Nieuniknioność typów głównych w teorii zupełnej]
\label{lem:mod-ott-principal-unavoidable}
Jeżeli $T$ jest zupełna, a $p$ jest głównym typem
częściowym nad $T$, to żaden model $T$ nie pomija $p$.
\end{lemat}
\begin{proof}
Niech $\theta$ świadczy o główności $p$.
Zdanie $\exists\bar x\,\theta$ zachodzi w pewnym
modelu $T$, więc przez zupełność teorii zachodzi
w każdym jej modelu. Dowolny jego świadek
realizuje wszystkie formuły $p$ z definicji
główności.
\end{proof}

Dla teorii niezupełnej wniosek ten jest fałszywy:
przykład \ref{ex:mod-isolated-incomplete} podaje
izolowany typ struktury jednoelementowej, pomijany
w modelu dwuelementowym tej samej teorii.

\section{Lemat o unikaniu typu przy skończonych wymaganiach}
\begin{lemat}[Unikanie typu w języku z nowymi stałymi]
\label{lem:mod-ott-avoidance}
Niech $C$ będzie zbiorem nowych stałych,
rozłącznym z językiem $\mathcal L$.
Niech $\Delta$ będzie skończonym zbiorem zdań
języka $\mathcal L(C)$ takim, że $T\cup\Delta$
jest spełnialny, a $p\subseteq F_n(\mathcal L)$
niegłównym typem nad $T$. Dla każdej krotki
$\bar t=(t_1,\ldots,t_n)$ termów domkniętych
języka $\mathcal L(C)$ istnieje $\varphi\in p$
taka, że
\[
T\cup\Delta\cup\{\neg\varphi(\bar t)\}
\]
jest spełnialny.
\end{lemat}
\begin{proof}
Przypuśćmy przeciwnie, że dla każdej $\varphi\in p$
zachodzi $T\cup\Delta\models\varphi(\bar t)$.
Niech $\bar c$ będzie listą wszystkich różnych nowych
stałych występujących w $\Delta$ lub w termach $\bar t$.
Jest to lista skończona, ponieważ $\Delta$ jest
skończony, a formuły i termy są skończonymi wyrażeniami.
Po ewentualnej zmianie nazw zmiennych związanych
zastępujemy te stałe parami różnymi świeżymi
zmiennymi $\bar z$, rozłącznymi z $\bar x$.
Oznaczmy przez $\delta(\bar z)$ wynik tej zamiany
w koniunkcji zdań $\Delta$, a przez
$s_i(\bar z)$ wynik zamiany w $t_i$.
Dla pustej koniunkcji używamy $\top$,
zgodnie z konwencją z rozdziału \ref{ch:mod-qe}.
W języku pierwotnym tworzymy formułę
\[
\theta(\bar x)=
\exists\bar z\,\left(
\delta(\bar z)\land\bigwedge_{i=1}^{n} x_i=s_i(\bar z)
\right).
\]
Zmienne $\bar x$ wybieramy tak, aby żadna z nich
nie występowała jako związana w $\delta$.
Jeżeli lista $\bar c$ jest pusta, nie dodajemy
kwantyfikatorów $\exists\bar z$.

Spełnialność $T\cup\Delta$ daje model, w którym
wartości nowych stałych świadczą o prawdziwości
$\theta$ przy krotce wartości termów $\bar t$.
Zatem $\theta$ jest realizowalna w modelu $T$.
Niech teraz dowolny model $\mathfrak A\models T$
i krotka $\bar a$ spełniają $\theta$.
Wybieramy świadków $\bar b$ dla $\bar z$.
Rozwijamy $\mathfrak A$ do języka $\mathcal L(C)$,
interpretując stałe $\bar c$ przez $\bar b$,
a pozostałe nowe stałe przez dowolny element
niepustego uniwersum $A$. Z semantyki podstawienia
(twierdzenie \ref{thm:sem-substitution}) oraz
indukcji po budowie termów wynika, że rozwinięcie
spełnia $\Delta$, a wartości termów $\bar t$
są równe odpowiednio współrzędnym $\bar a$.
Przy podstawianiu termów za zmienne w $\varphi$
wykorzystujemy to, że termy $\bar t$ są domknięte;
nie zawierają zmiennych, które mogłyby zostać
przechwycone przez kwantyfikator.

Z założenia rozwinięcie spełnia $\varphi(\bar t)$
dla każdej $\varphi\in p$. Zatem
$\mathfrak A\models\varphi(\bar a)$ dla każdej
takiej $\varphi$. Dowolność modelu i krotki daje
\[
T\models\forall\bar x\,(\theta(\bar x)\to\varphi(\bar x))
\qquad(\varphi\in p).
\]
Realizowalna formuła $\theta$ świadczyłaby więc
o główności $p$, wbrew założeniu.
\end{proof}

Ważne jest to, że $\Delta$ jest skończony, natomiast
$T$ może być nieskończony. Wszystkie wymagania z $T$
są zachowywane w każdym kroku. W lemacie nie
zakładamy ani przeliczalności języka, ani zupełności
teorii; przeliczalność będzie potrzebna przy
organizacji kolejnych etapów konstrukcji.

\section{Twierdzenie o pomijaniu typów}
\begin{twierdzenie}[Pomijanie przeliczalnej rodziny typów]
\label{thm:mod-ott}
Niech $\mathcal L$ będzie co najwyżej przeliczalnym
językiem, a $T$ spełnialną teorią tego języka.
Niech $(p_j)_{j\in J}$ będzie co najwyżej przeliczalną
rodziną niegłównych typów częściowych nad $T$,
gdzie $p_j\subseteq F_{n_j}(\mathcal L)$ i $n_j\geq1$.
Istnieje model $\mathfrak M\models T$ o mocy
co najwyżej $\aleph_0$, który pomija każdy $p_j$.
Jeżeli $T$ nie ma modeli skończonych, to
$|M|=\aleph_0$.
\end{twierdzenie}
\begin{proof}
Dodajemy przeliczalnie nieskończony zbiór nowych
stałych $C=\{c_0,c_1,\ldots\}$.
Rozszerzony język $\mathcal L(C)$ jest nadal
przeliczalny. Jego termy i formuły tworzą zbiory
co najwyżej przeliczalne, jako skończone ciągi
symboli z przeliczalnego alfabetu (por. lemat
\ref{lem:ls-formula-count}).

Ustalamy wyliczenie trzech rodzajów wymagań:
\begin{enumerate}[label=(\roman*)]
\item rozstrzygnąć każde zdanie $\sigma$ języka
$\mathcal L(C)$;
\item zapewnić świadka dla każdego zdania postaci
$\exists x\,\psi(x)$ języka $\mathcal L(C)$,
gdzie jedyną możliwą zmienną wolną $\psi$ jest $x$;
\item dla każdego $j\in J$ i każdej krotki $\bar t$
długości $n_j$ złożonej z termów domkniętych
języka $\mathcal L(C)$ wymusić
$\neg\varphi(\bar t)$ dla pewnej $\varphi\in p_j$.
\end{enumerate}
Zbiór wszystkich wymagań jest przeliczalny:
w trzecim przypadku używamy przeliczalności
rodziny $J$ i skończoności $n_j$.
Wyliczenie może zawierać powtórzenia.
Budujemy ciąg skończonych zbiorów zdań
\[
\varnothing=\Delta_0\subseteq\Delta_1
\subseteq\Delta_2\subseteq\cdots
\]
tak, aby $T\cup\Delta_k$ był spełnialny dla każdego
$k$. Na etapie $k$ obsługujemy $k$-te wymaganie.

Dla wymagania (i) wybieramy jedno ze zdań
$\sigma,\neg\sigma$, którego dodanie zachowuje
spełnialność. Co najmniej jedno jest dopuszczalne,
ponieważ dowolny model $T\cup\Delta_k$ spełnia
dokładnie jedno z nich.

Dla wymagania (ii) wybieramy stałą $c\in C$,
która nie występuje w $\Delta_k$ ani w zdaniu
$\exists x\,\psi(x)$. Jest to możliwe, ponieważ
te wyrażenia używają tylko skończenie wielu
nowych stałych, a $T$ nie używa żadnej z nich.
Dodajemy zdanie Henkina
\[
(\exists x\,\psi(x))\to\psi(c).
\]
Aby sprawdzić spełnialność, wybieramy dowolny
model $T\cup\Delta_k$. Jeżeli przesłanka jest
prawdziwa, zmieniamy interpretację $c$ na
wartość świadka. Jeżeli jest fałszywa, nie musimy
nic zmieniać. W obu przypadkach wcześniejsze
wymagania pozostają prawdziwe.
Jest to semantyczna wersja kroku z lematu
\ref{lem:henkin-one-witness}.

Dla wymagania (iii) stosujemy lemat
\ref{lem:mod-ott-avoidance} do $p_j$,
$\Delta_k$ i krotki $\bar t$. Wybieramy formułę
$\varphi\in p_j$ wskazaną przez lemat i dodajemy
$\neg\varphi(\bar t)$.
Możemy ustalić dobre uporządkowanie formuł
rozszerzonego języka i zawsze brać pierwszą
dopuszczalną formułę. W ten sposób wybory
dokonywane są w obrębie zbioru.

Niech
\[
\Gamma=T\cup\bigcup_{k<\omega}\Delta_k.
\]
Każdy skończony fragment $\Gamma$ należy do
$T\cup\Delta_k$ dla dostatecznie dużego $k$.
Z twierdzenia o zwartości \ref{thm:compactness}
istnieje model $\mathfrak D\models\Gamma$
w języku rozszerzonym. Nie możemy po prostu
wziąć reduktu $\mathfrak D$: mogą w nim istnieć
elementy niebędące wartościami żadnych termów
domkniętych, których wymagania (iii) nie dotyczą.

Niech $D_0$ będzie zbiorem wartości wszystkich
termów domkniętych w $\mathfrak D$.
Jest niepusty, gdyż zawiera wartość $c_0$,
i co najwyżej przeliczalny. Zawiera wartości
wszystkich stałych i jest zamknięty na działania:
jeżeli $a_i=t_i^{\mathfrak D}$, to
$f^{\mathfrak D}(\bar a)$ jest wartością termu
$f(t_1,\ldots,t_m)$. Definiuje więc podstrukturę
$\mathfrak D_0\subseteq\mathfrak D$.

Pokażemy, że $\mathfrak D_0\preccurlyeq\mathfrak D$
za pomocą kryterium Tarskiego--Vaughta
\ref{thm:model-tarski-vaught}.
Niech $\mathfrak D\models\exists x\,\eta(x,\bar a)$,
gdzie $\bar a$ jest krotką z $D_0$. Każda współrzędna
jest wartością pewnego termu domkniętego $t_i$.
Podstawiając te termy za odpowiednie zmienne
w $\eta$, otrzymujemy formułę $\psi(x)$, której
jedyną możliwą zmienną wolną jest $x$.
Jeżeli zachodzi potrzeba, wcześniej przemianowujemy
zmienne związane. Semantyka podstawienia daje
$\mathfrak D\models\exists x\,\psi(x)$.
Wymaganie (ii) dla tego zdania zapewniło
aksjomat $(\exists x\,\psi(x))\to\psi(c)$
należący do $\Gamma$. Zatem
$\mathfrak D\models\psi(c)$, czyli
$\mathfrak D\models\eta(c^{\mathfrak D},\bar a)$.
Ponieważ $c^{\mathfrak D}\in D_0$, otrzymaliśmy
świadka w $D_0$ i kryterium daje elementarność.

Niech $\mathfrak M$ będzie reduktem $\mathfrak D_0$
do języka $\mathcal L$. Przez elementarność
$\mathfrak D_0\models T$, więc $\mathfrak M\models T$.
Ustalmy $j\in J$ i dowolną krotkę $\bar a\in M^{n_j}$.
Wybieramy krotkę termów domkniętych $\bar t$
o wartościach $\bar a$ w $\mathfrak D$.
Wymaganie (iii) zapewniło zdanie
$\neg\varphi(\bar t)\in\Gamma$
dla pewnej $\varphi\in p_j$. W konsekwencji
$\mathfrak D\models\neg\varphi(\bar t)$.
Z elementarności $\mathfrak D_0$ i semantyki
podstawienia wynika $\mathfrak M\models\neg\varphi(\bar a)$.
Krotka $\bar a$ nie realizuje $p_j$.
Dowolność $\bar a,j$ dowodzi pomijania całej
rodziny. Ograniczenie mocy wynika z konstrukcji $D_0$;
jeżeli $T$ nie ma skończonych modeli, model
musi być przeliczalnie nieskończony.
\end{proof}

Konstrukcja jest wariantem metody Henkina
z rozdziału \ref{ch:completeness}. Zamiast
konstruować model jako iloraz termów względem
równości dowodliwej, korzystamy z już udowodnionej
zwartości i bierzemy elementarną podstrukturę
generowaną przez termy domknięte. Warunek
rozstrzygania zdań zapewnia dodatkowo, że $\Gamma$
jest zupełna semantycznie; nie jest on niezbędny
do ostatniego argumentu o podstrukturze.
Kluczowe są wymagania Henkina oraz wymagania
unikające realizacji typów.

\begin{figure}[htbp]
\centering
\begin{tikzpicture}[>=Stealth,x=1cm,y=1cm]
\foreach \x/\i in {0/0,2.7/1,5.4/2,8.1/3}{
 \node[draw,rounded corners,minimum width=1.3cm]
 (d\i) at (\x,0) {$\Delta_{\i}$};
}
\draw[->,thick] (d0)--node[above,font=\small] {zdanie}(d1);
\draw[->,thick] (d1)--node[above,font=\small] {świadek}(d2);
\draw[->,thick] (d2)--node[above,font=\small] {unikanie}(d3);
\node at (9.6,0) {$\cdots$};
\draw[->,dashed,thick] (d3)--(9.2,0);
\node[align=center,text width=9cm,font=\small] at (4.8,-1)
 {Każdy etap zachowuje spełnialność $T\cup\Delta_k$.};
\end{tikzpicture}
\caption{Trzy rodzaje wymagań w konstrukcji modelu
pomijającego typy. Kolejność na rysunku jest przykładowa:
wyliczenie musi objąć wszystkie wymagania.}
\label{fig:mod-ott-henkin}
\end{figure}

\begin{wniosek}[Charakterystyka typów pomijalnych]
\label{cor:mod-ott-complete}
Niech $T$ będzie zupełną teorią w co najwyżej
przeliczalnym języku. Typ częściowy nad $T$
można pominąć w pewnym modelu $T$ dokładnie wtedy,
gdy jest niegłówny. Można wtedy wybrać model
o mocy co najwyżej $\aleph_0$.
Dla typu zupełnego oznacza to dokładnie
nieizolowanie w przestrzeni typów.
\end{wniosek}
\begin{proof}
Typ główny jest realizowany w każdym modelu $T$
na mocy lematu \ref{lem:mod-ott-principal-unavoidable}.
Typ niegłówny można pominąć na mocy twierdzenia
\ref{thm:mod-ott}, zastosowanego do rodziny
jednoelementowej. Dla typu zupełnego stosujemy
lemat \ref{lem:mod-ott-interior}.
\end{proof}

\section{Przykłady i ograniczenia twierdzenia}
\begin{przyklad}[Typ elementu spoza standardowej arytmetyki]
\label{ex:mod-ott-arithmetic}
Niech $T=\operatorname{Th}(\mathbb N,0,S,+,\cdot)$,
gdzie $S(n)=n+1$, a $\underline n=S^n(0)$.
Zbiór
\[
p(x)=\{x\ne\underline n:n\in\mathbb N\}
\]
jest typem częściowym nad $T$:
każdy skończony podzbiór wymaga uniknięcia
skończenie wielu liczb naturalnych, co można
zrealizować w modelu standardowym. Lemat
\ref{lem:mod-type-compactness} daje realizowalność
całego typu w pewnym modelu $T$.
Model standardowy pomija jednak $p$, bo każdy
jego element jest wartością pewnego numeralu.
Z lematu \ref{lem:mod-ott-principal-unavoidable}
wynika, że $p$ jest niegłówny. Z twierdzenia
\ref{thm:mod-ott} istnieje model przeliczalny
pomijający $p$ (w tym przypadku znamy go wprost).
Każdy model $T$ pomijający $p$ jest izomorficzny
z modelem standardowym: jego uniwersum składa
się z wartości numerali, różniących się parami,
a wszystkie równania opisujące działania $S,+,\cdot$
na numeralach są konsekwencjami $T$.
Typ $p$ ma także realizację w przeliczalnym
modelu niestandardowym, przez zwartość i dolne
twierdzenie Löwenheima--Skolema
\ref{thm:downward-ls} zastosowane do modelu
zawierającego realizację typu.
\end{przyklad}

\begin{przyklad}[Nie można pomijać dowolnie wielu typów]
\label{ex:mod-ott-uncountable-family}
Niech $\mathcal L=\{P_i:i\in\mathbb N\}$ będzie
językiem z przeliczalną rodziną jednoargumentowych
symboli relacyjnych. Rozważmy strukturę
$\mathfrak A$ o uniwersum wszystkich skończonych
podzbiorów $\mathbb N$ z interpretacją
\[
P_i^{\mathfrak A}(F)\quad\Longleftrightarrow\quad i\in F.
\]
Niech $T=\operatorname{Th}(\mathfrak A)$.
Dla każdego $D\subseteq\mathbb N$ tworzymy typ częściowy
\[
p_D(x)=\{P_i(x):i\in D\}
\cup\{\neg P_i(x):i\notin D\}.
\]
Jest realizowalny w pewnym modelu $T$, gdyż
każdy skończony fragment można zrealizować
w $\mathfrak A$ przez skończony zbiór dodatnich
wymagań. Stosujemy lemat \ref{lem:mod-type-compactness}.

Wykażemy niegłówność każdego $p_D$.
Niech $\theta(x)$ będzie dowolną formułą
realizowalną w pewnym modelu $T$.
Z zupełności $T$ istnieje skończony zbiór
$F\subseteq\mathbb N$ spełniający $\theta$ w
$\mathfrak A$. Niech $I$ będzie skończonym
zbiorem indeksów symboli $P_i$ występujących
w $\theta$. Wybieramy $k\notin I$. Odwzorowanie
$G\mapsto G\mathbin{\triangle}\{k\}$, gdzie
$\triangle$ oznacza różnicę symetryczną, jest
bijekcją uniwersum $\mathfrak A$ i automorfizmem
reduktu do symboli z $I$: zachowuje przynależność
do każdego $P_i$ dla $i\in I$.
Z niezmienniczości względem izomorfizmu
(twierdzenie \ref{thm:model-iso-invariance})
oba zbiory $F$ i $F\mathbin{\triangle}\{k\}$
spełniają $\theta$. Mają przeciwne wartości $P_k$,
więc co najmniej jeden nie realizuje $p_D$.
W konsekwencji $\theta$ nie wymusza wszystkich
formuł $p_D$. Dowolność $\theta$ daje niegłówność.

Rodzina $(p_D)_{D\subseteq\mathbb N}$ ma moc
$2^{\aleph_0}$. Żaden model $T$ nie pomija wszystkich
tych typów: dla dowolnego elementu $a$ modelu
$\mathfrak M\models T$ wystarczy położyć
\[
D_a=\{i\in\mathbb N:\mathfrak M\models P_i(a)\}.
\]
Element $a$ realizuje $p_{D_a}$ z definicji.
Twierdzenia \ref{thm:mod-ott} nie można więc
rozszerzyć bezwarunkowo na dowolne nieprzeliczalne
rodziny typów, nawet w przeliczalnym języku
i dla teorii zupełnej.
\end{przyklad}

W dalszej części rozdziału zastosujemy twierdzenie
o pomijaniu typów do charakterystyki modeli atomowych
i pierwszych.
